\documentclass{beamer}

\usetheme{Madrid}
\usecolortheme{seahorse}

\usepackage[utf8]{inputenc}
\usepackage[spanish]{babel}
\usepackage{amsmath}
\usepackage{multicol}

\title{Sistemas de Ecuaciones Lineales en Agronomía}
\subtitle{Métodos algebraicos y análisis de soluciones}
\author{Matemática Aplicada}
\date{}

\begin{document}
	
	\frame{\titlepage}
	
	%------------------------------------------------
	\begin{frame}{Objetivo de la clase}
		Analizar y resolver situaciones agronómicas mediante sistemas de ecuaciones lineales con dos incógnitas.
		
		\bigskip
		
		Al finalizar la clase, el estudiante será capaz de:
		
		\begin{itemize}
			\item Identificar sistemas independientes, dependientes e inconsistentes.
			\item Resolver sistemas mediante sustitución, igualación y reducción.
			\item Interpretar soluciones en contextos agronómicos.
			\item Detectar patrones, relaciones e inconsistencias.
		\end{itemize}
	\end{frame}
	
	%------------------------------------------------
	\begin{frame}{Situación inicial: mezcla de fertilizantes}
		Un productor necesita preparar una mezcla de fertilizantes para un cultivo de tomates.
		
		\bigskip
		
		\begin{itemize}
			\item Fertilizante A: aporta nitrógeno y fósforo.
			\item Fertilizante B: aporta nutrientes en distinta proporción.
			\item Se busca cumplir una dosis exacta para evitar déficit o exceso.
		\end{itemize}
		
		\bigskip
		
		\textbf{Pregunta:}  
		¿Cómo determinar cuántos kilogramos de cada fertilizante usar?
	\end{frame}
	
	%------------------------------------------------
	\begin{frame}{Sistema de ecuaciones lineales}
		Un sistema con dos incógnitas se expresa como:
		
		\[
		\begin{cases}
			ax + by = c \\
			dx + ey = f
		\end{cases}
		\]
		
		\bigskip
		
		En agronomía:
		
		\begin{itemize}
			\item $x$: cantidad del insumo A.
			\item $y$: cantidad del insumo B.
			\item Las ecuaciones representan restricciones de producción, costos, nutrientes o superficie.
		\end{itemize}
	\end{frame}
	
	%------------------------------------------------
	\begin{frame}{Tipos de sistemas}
		\begin{block}{Sistema independiente}
			Tiene una única solución.  
			Las rectas se intersectan en un punto.
		\end{block}
		
		\begin{block}{Sistema dependiente}
			Tiene infinitas soluciones.  
			Las ecuaciones representan la misma recta.
		\end{block}
		
		\begin{block}{Sistema inconsistente}
			No tiene solución.  
			Las rectas son paralelas.
		\end{block}
	\end{frame}
	
	%------------------------------------------------
	\begin{frame}{Ejemplo 1: sistema independiente}
		Un agrónomo mezcla dos fertilizantes.
		
		\begin{itemize}
			\item $x$: kg de fertilizante A.
			\item $y$: kg de fertilizante B.
		\end{itemize}
		
		\[
		\begin{cases}
			x + y = 10 \\
			2x + y = 14
		\end{cases}
		\]
		
		Restando la primera ecuación de la segunda:
		
		\[
		x = 4
		\]
		
		Luego:
		
		\[
		4 + y = 10
		\]
		
		\[
		y = 6
		\]
		
		\bigskip
		
		\textbf{Interpretación:} se deben usar 4 kg de A y 6 kg de B.
	\end{frame}
	
	%------------------------------------------------
	\begin{frame}{Ejemplo 2: sistema dependiente}
		Una finca registra las siguientes condiciones de mezcla:
		
		\[
		\begin{cases}
			x + y = 8 \\
			2x + 2y = 16
		\end{cases}
		\]
		
		La segunda ecuación es el doble de la primera:
		
		\[
		2(x+y)=16
		\]
		
		\[
		x+y=8
		\]
		
		\bigskip
		
		\textbf{Interpretación:} existen infinitas combinaciones posibles de insumos que cumplen la condición.
	\end{frame}
	
	%------------------------------------------------
	\begin{frame}{Ejemplo 3: sistema inconsistente}
		Un técnico registra las siguientes restricciones:
		
		\[
		\begin{cases}
			x + y = 5 \\
			x + y = 9
		\end{cases}
		\]
		
		\bigskip
		
		Esto indica que la misma suma debería ser 5 y 9 al mismo tiempo.
		
		\bigskip
		
		\textbf{Interpretación:} los datos son contradictorios. No existe una combinación posible.
	\end{frame}
	
	%------------------------------------------------
	\begin{frame}{Método de sustitución}
		\textbf{Problema:}  
		En una parcela se siembran maíz y trigo.
		
		\begin{itemize}
			\item $x$: kg de semilla de maíz.
			\item $y$: kg de semilla de trigo.
		\end{itemize}
		
		Se sembraron 50 kg en total. Además, se usaron 10 kg más de maíz que de trigo.
		
		\[
		\begin{cases}
			x + y = 50 \\
			x = y + 10
		\end{cases}
		\]
		
		Sustituyendo:
		
		\[
		(y+10)+y=50
		\]
		
		\[
		2y=40
		\]
		
		\[
		y=20
		\]
		
		\[
		x=30
		\]
		
		\textbf{Respuesta:} 30 kg de maíz y 20 kg de trigo.
	\end{frame}
	
	%------------------------------------------------
	\begin{frame}{Ejercicio 1: sustitución}
		En un vivero se producen plantas de tomate y lechuga.
		
		\begin{itemize}
			\item $x$: bandejas de tomate.
			\item $y$: bandejas de lechuga.
		\end{itemize}
		
		En total se produjeron 80 bandejas.  
		La cantidad de bandejas de tomate fue 20 más que la de lechuga.
		
		\[
		\begin{cases}
			x + y = 80 \\
			x = y + 20
		\end{cases}
		\]
		
		\bigskip
		
		\textbf{Tarea:} resolver por sustitución e interpretar el resultado.
	\end{frame}
	
	%------------------------------------------------
	\begin{frame}{Solución ejercicio 1}
		\[
		x = y + 20
		\]
		
		Sustituimos:
		
		\[
		(y+20)+y=80
		\]
		
		\[
		2y=60
		\]
		
		\[
		y=30
		\]
		
		\[
		x=50
		\]
		
		\bigskip
		
		\textbf{Interpretación:} se produjeron 50 bandejas de tomate y 30 de lechuga.
	\end{frame}
	
	%------------------------------------------------
	\begin{frame}{Método de igualación}
		\textbf{Problema:}  
		Dos expresiones permiten estimar la cantidad de agua disponible para riego.
		
		\[
		\begin{cases}
			x = 30 - y \\
			x = 2y + 6
		\end{cases}
		\]
		
		Igualamos:
		
		\[
		30-y=2y+6
		\]
		
		\[
		24=3y
		\]
		
		\[
		y=8
		\]
		
		\[
		x=22
		\]
		
		\bigskip
		
		\textbf{Interpretación:} el punto común del modelo es $x=22$, $y=8$.
	\end{frame}
	
	%------------------------------------------------
	\begin{frame}{Ejercicio 2: igualación}
		En una explotación agrícola se comparan dos modelos de costos.
		
		\[
		\begin{cases}
			x = 100 - 2y \\
			x = 3y + 20
		\end{cases}
		\]
		
		\bigskip
		
		\textbf{Tarea:}
		
		\begin{itemize}
			\item Resolver por igualación.
			\item Determinar si el sistema es independiente.
			\item Interpretar el resultado en términos de costos y producción.
		\end{itemize}
	\end{frame}
	
	%------------------------------------------------
	\begin{frame}{Solución ejercicio 2}
		Igualamos:
		
		\[
		100-2y=3y+20
		\]
		
		\[
		80=5y
		\]
		
		\[
		y=16
		\]
		
		Sustituimos:
		
		\[
		x=3(16)+20
		\]
		
		\[
		x=68
		\]
		
		\bigskip
		
		\textbf{Solución:}
		
		\[
		(x,y)=(68,16)
		\]
		
		\textbf{Tipo de sistema:} independiente.
	\end{frame}
	
	%------------------------------------------------
	\begin{frame}{Método de reducción}
		\textbf{Problema:}  
		Una muestra de suelo requiere dos nutrientes:
		
		\[
		\begin{cases}
			2x + y = 11 \\
			x - y = 1
		\end{cases}
		\]
		
		Sumamos ambas ecuaciones:
		
		\[
		3x=12
		\]
		
		\[
		x=4
		\]
		
		Sustituimos:
		
		\[
		4-y=1
		\]
		
		\[
		y=3
		\]
		
		\bigskip
		
		\textbf{Interpretación:} se requieren 4 unidades del primer nutriente y 3 del segundo.
	\end{frame}
	
	%------------------------------------------------
	\begin{frame}{Ejercicio 3: reducción}
		Para fertilizar una hectárea se combinan dos productos.
		
		\begin{itemize}
			\item $x$: litros del producto A.
			\item $y$: litros del producto B.
		\end{itemize}
		
		El técnico obtiene el sistema:
		
		\[
		\begin{cases}
			3x + 2y = 24 \\
			x - 2y = 0
		\end{cases}
		\]
		
		\bigskip
		
		\textbf{Tarea:} resolver por reducción e interpretar la solución.
	\end{frame}
	
	%------------------------------------------------
	\begin{frame}{Solución ejercicio 3}
		Sistema:
		
		\[
		\begin{cases}
			3x + 2y = 24 \\
			x - 2y = 0
		\end{cases}
		\]
		
		Sumamos:
		
		\[
		4x=24
		\]
		
		\[
		x=6
		\]
		
		Sustituimos:
		
		\[
		6-2y=0
		\]
		
		\[
		2y=6
		\]
		
		\[
		y=3
		\]
		
		\bigskip
		
		\textbf{Respuesta:} 6 litros del producto A y 3 litros del producto B.
	\end{frame}
	
	%------------------------------------------------
	\begin{frame}{Actividad integradora}
		Una finca cosecha papa y cebolla.
		
		\begin{itemize}
			\item $x$: sacos de papa.
			\item $y$: sacos de cebolla.
		\end{itemize}
		
		En total se cosecharon 120 sacos.
		
		Cada saco de papa genera \$8.000 y cada saco de cebolla genera \$12.000.
		
		La ganancia total fue de \$1.200.000.
		
		\[
		\begin{cases}
			x+y=120 \\
			8000x+12000y=1200000
		\end{cases}
		\]
		
		\bigskip
		
		\textbf{Resolver e interpretar.}
	\end{frame}
	
	%------------------------------------------------
	\begin{frame}{Solución actividad integradora}
		Dividimos la segunda ecuación por 4000:
		
		\[
		2x+3y=300
		\]
		
		Sistema equivalente:
		
		\[
		\begin{cases}
			x+y=120 \\
			2x+3y=300
		\end{cases}
		\]
		
		Multiplicamos la primera por 2:
		
		\[
		2x+2y=240
		\]
		
		Restamos:
		
		\[
		y=60
		\]
		
		\[
		x=60
		\]
		
		\bigskip
		
		\textbf{Interpretación:} se cosecharon 60 sacos de papa y 60 de cebolla.
	\end{frame}
	
	%------------------------------------------------
	\begin{frame}{Análisis de patrones e inconsistencias}
		Al resolver sistemas, debemos observar:
		
		\begin{itemize}
			\item Si aparece una única solución: sistema independiente.
			\item Si una ecuación es múltiplo de otra: sistema dependiente.
			\item Si aparece una contradicción como $0=5$: sistema inconsistente.
			\item Si el resultado no tiene sentido en el contexto, debe revisarse el modelo.
		\end{itemize}
		
		\bigskip
		
		\textbf{Ejemplo de inconsistencia contextual:}
		
		\[
		x=-4
		\]
		
		No puede representar una cantidad negativa de fertilizante o sacos cosechados.
	\end{frame}
	
	%------------------------------------------------
	\begin{frame}{Ejercicio 4: detectar el tipo de sistema}
		Clasifica cada sistema según corresponda:
		
		\[
		A)
		\begin{cases}
			x+y=10 \\
			2x+2y=20
		\end{cases}
		\]
		
		\[
		B)
		\begin{cases}
			x+y=10 \\
			x+y=15
		\end{cases}
		\]
		
		\[
		C)
		\begin{cases}
			x+y=10 \\
			x-y=2
		\end{cases}
		\]
		
		\bigskip
		
		\textbf{Opciones:} independiente, dependiente o inconsistente.
	\end{frame}
	
	%------------------------------------------------
	\begin{frame}{Solución ejercicio 4}
		\[
		A)
		\begin{cases}
			x+y=10 \\
			2x+2y=20
		\end{cases}
		\]
		
		Sistema dependiente.
		
		\bigskip
		
		\[
		B)
		\begin{cases}
			x+y=10 \\
			x+y=15
		\end{cases}
		\]
		
		Sistema inconsistente.
		
		\bigskip
		
		\[
		C)
		\begin{cases}
			x+y=10 \\
			x-y=2
		\end{cases}
		\]
		
		Sistema independiente.
	\end{frame}
	
	%------------------------------------------------
	\begin{frame}{Cierre de la clase}
		\textbf{Preguntas de reflexión:}
		
		\begin{itemize}
			\item ¿Qué método fue más eficiente en cada situación?
			\item ¿Cómo se interpreta una solución en agronomía?
			\item ¿Qué riesgos existen si se trabaja con datos inconsistentes?
			\item ¿Por qué no siempre una solución algebraica es válida en el contexto real?
		\end{itemize}
	\end{frame}
	
	%------------------------------------------------
	\begin{frame}{Evaluación formativa}
		\begin{center}
			\begin{tabular}{|p{6cm}|c|c|}
				\hline
				\textbf{Criterio} & \textbf{Logrado} & \textbf{Por mejorar} \\
				\hline
				Identifica el tipo de sistema & & \\
				\hline
				Resuelve por sustitución & & \\
				\hline
				Resuelve por igualación & & \\
				\hline
				Resuelve por reducción & & \\
				\hline
				Interpreta resultados agronómicos & & \\
				\hline
				Detecta inconsistencias & & \\
				\hline
			\end{tabular}
		\end{center}
	\end{frame}
	
	%------------------------------------------------
	\begin{frame}{Tarea}
		Plantea un problema agronómico real que pueda resolverse mediante un sistema de ecuaciones lineales.
		
		\bigskip
		
		Debe incluir:
		
		\begin{itemize}
			\item Contexto agrícola.
			\item Definición de variables.
			\item Sistema de ecuaciones.
			\item Resolución algebraica.
			\item Interpretación del resultado.
			\item Clasificación del sistema.
		\end{itemize}
	\end{frame}
	
\end{document}